\documentclass[11pt,a4paper]{article}

% ---- Packages ----
\usepackage[utf8]{inputenc}
\usepackage[T1]{fontenc}
\usepackage{amsmath,amssymb,amsthm}
\usepackage{mathtools}
\usepackage[margin=1in,headheight=14pt]{geometry}
\usepackage{hyperref}
\usepackage{enumitem}
\usepackage{xcolor}
\usepackage{tcolorbox}
\usepackage{fancyhdr}
\usepackage{booktabs}

% ---- Hyperref ----
\hypersetup{
    colorlinks=true,
    linkcolor=red,
    citecolor=blue,
    urlcolor=blue
}

% ---- Header ----
\pagestyle{fancy}
\fancyhf{}
\fancyhead[L]{\textit{Lesson 10: Stochastic Approximation}}
\fancyhead[R]{\thepage}
\renewcommand{\headrulewidth}{0.4pt}

% ---- Theorem Environments ----
\theoremstyle{definition}
\newtheorem{definition}{Definition}[section]
\newtheorem{example}{Example}[section]
\newtheorem{exercise}{Exercise}[section]

\theoremstyle{plain}
\newtheorem{theorem}{Theorem}[section]
\newtheorem{lemma}[theorem]{Lemma}
\newtheorem{proposition}[theorem]{Proposition}
\newtheorem{corollary}[theorem]{Corollary}

\theoremstyle{remark}
\newtheorem{remark}{Remark}[section]

% ---- Colored Boxes ----
\tcbuselibrary{theorems,skins,breakable}

\newtcolorbox{keyresult}[1][]{
    colback=green!5!white,
    colframe=green!50!black,
    fonttitle=\bfseries,
    title=#1,
    breakable
}

\newtcolorbox{rlconnection}[1][]{
    colback=blue!5!white,
    colframe=blue!50!black,
    fonttitle=\bfseries,
    title=#1,
    breakable
}

\newtcolorbox{warning}[1][]{
    colback=red!5!white,
    colframe=red!50!black,
    fonttitle=\bfseries,
    title=#1,
    breakable
}

% ---- Operators ----
\newcommand{\R}{\mathbb{R}}
\newcommand{\N}{\mathbb{N}}
\newcommand{\Q}{\mathbb{Q}}
\newcommand{\Z}{\mathbb{Z}}
\newcommand{\C}{\mathbb{C}}
\newcommand{\Prob}{\mathbb{P}}
\newcommand{\E}{\mathbb{E}}
\newcommand{\Var}{\mathrm{Var}}
\newcommand{\indicator}{\mathbf{1}}
\newcommand{\cas}{\xrightarrow{\text{a.s.}}}
\newcommand{\cprob}{\xrightarrow{\Prob}}
\newcommand{\cdist}{\xrightarrow{d}}

% ---- Title ----
\title{\textbf{Lesson 10: Stochastic Approximation}\\[10pt]
\large Robbins--Monro, the ODE Method, and Two-Timescale Analysis\\[5pt]
\normalsize Session 10 of 102 $\mid$ Phase I: Mathematical Foundations $\mid$ 8\,h\\
\normalsize Primary text: Borkar, \emph{Stochastic Approximation: A Dynamical Systems Viewpoint}, Ch.\,1--3, 6}
\author{Curriculum: A Rigorous Learning Path to Multi-Agent Deep RL}
\date{}

\begin{document}

\maketitle

\begin{abstract}
Stochastic approximation (SA) is the mathematical engine behind virtually every
reinforcement learning algorithm. Temporal-Difference learning, Q-learning, SARSA,
actor-critic methods, and policy gradient algorithms are all instances of stochastic
approximation schemes. This document develops the theory of SA from the classical
Robbins--Monro algorithm through the ODE method of Ljung and Kushner--Clark, convergence
rate analysis via the central limit theorem for SA, and Borkar's two-timescale framework
that underlies actor-critic architectures. We provide complete proofs of all main
convergence theorems under precise assumptions, emphasizing the dynamical systems
viewpoint: a stochastic approximation iteration with decreasing step sizes asymptotically
tracks the solutions of an associated ordinary differential equation. Every major result
is connected to specific RL algorithms through detailed worked examples.
\end{abstract}

\tableofcontents
\newpage

%% ============================================================
%% SECTION 1: INTRODUCTION
%% ============================================================
\section{Introduction}
\label{sec:intro}

The central challenge in reinforcement learning is optimization under uncertainty: an
agent must find an optimal policy when the environment dynamics are unknown and the only
information available comes through noisy, sequential interactions. This is precisely the
setting of \emph{stochastic approximation} --- finding the root (or optimum) of a
function that can only be observed through noisy samples.

The field of stochastic approximation was born with the seminal 1951 paper of Robbins
and Monro, who proposed an iterative scheme for finding the root of a regression
function from noisy observations. This idea was extended to optimization by Kiefer and
Wolfowitz, and the theoretical foundations were deepened by Ljung, Kushner and Clark,
and Borkar, who established the powerful \emph{ODE method}: under appropriate conditions,
a stochastic approximation algorithm converges to the same limit as the solution of an
associated ordinary differential equation.

\medskip
\noindent\textbf{Why this matters for RL.} Consider the following RL update rules:
\begin{itemize}[leftmargin=2em]
\item \textbf{TD(0):} $V(s_n) \leftarrow V(s_n) + \alpha_n \bigl[R_{n+1} + \gamma V(s_{n+1}) - V(s_n)\bigr]$
\item \textbf{Q-learning:} $Q(s_n, a_n) \leftarrow Q(s_n,a_n) + \alpha_n \bigl[R_{n+1} + \gamma \max_{a'} Q(s_{n+1},a') - Q(s_n,a_n)\bigr]$
\item \textbf{Actor-Critic:} Critic updates value, actor updates policy, at \emph{different rates}
\item \textbf{REINFORCE:} $\theta_{n+1} = \theta_n + \alpha_n \widehat{\nabla J(\theta_n)}$
\end{itemize}
Every one of these is a stochastic approximation iteration. The convergence theorems
in this document therefore provide the theoretical foundation for understanding
\emph{when} and \emph{why} RL algorithms converge.

\medskip
\noindent\textbf{Document outline.}
Section~\ref{sec:robbins-monro} introduces the Robbins--Monro algorithm and proves its
convergence under classical conditions. Section~\ref{sec:ode-method} develops the ODE
method, the most powerful general-purpose tool for SA convergence analysis.
Section~\ref{sec:convergence-rate} establishes the convergence rate via the CLT for SA
and Polyak--Ruppert averaging. Section~\ref{sec:two-timescale} presents Borkar's
two-timescale SA theory, which is the backbone of actor-critic algorithms.
Section~\ref{sec:linear-sa} treats the important special case of linear SA.
Section~\ref{sec:rl-applications} maps the general theory to specific RL algorithms
in detail. Section~\ref{sec:exercises} provides exercises, and
Section~\ref{sec:summary} summarizes all key results.

\medskip
\noindent\textbf{Prerequisites.} Probability spaces, conditional expectation,
martingale difference sequences, the law of large numbers, and basic ODE theory
(existence and uniqueness of solutions, Lyapunov stability).

%% ============================================================
%% SECTION 2: THE ROBBINS-MONRO ALGORITHM
%% ============================================================
\section{The Robbins--Monro Algorithm}
\label{sec:robbins-monro}

\subsection{Problem formulation}

Suppose we wish to find a root of a function $h : \R^d \to \R^d$, i.e., a point
$x^* \in \R^d$ such that $h(x^*) = 0$. The difficulty is that we cannot evaluate $h$
directly. Instead, for each query point $x$, we observe a noisy sample
\[
  H(x, Y) \quad \text{where} \quad h(x) = \E[H(x, Y)]
\]
and $Y$ is a random variable (or random element) representing the noise source.

\begin{example}[Stochastic root-finding]
\label{ex:root-finding}
Let $h(\theta) = \E[f(X) - \theta]$ where $X$ is a random variable with unknown
distribution. Then $h(\theta^*) = 0$ when $\theta^* = \E[f(X)]$. Given i.i.d.\ samples
$X_1, X_2, \ldots$, the noisy observation is $H(\theta, X_n) = f(X_n) - \theta$.
\end{example}

\begin{definition}[Step-size sequence]
\label{def:step-size}
A \emph{step-size sequence} (or \emph{learning rate schedule}) is a sequence
$\{a_n\}_{n \geq 0}$ of positive real numbers satisfying:
\begin{enumerate}[label=(\roman*)]
\item $\displaystyle\sum_{n=0}^{\infty} a_n = \infty$ \quad (the steps are large enough to reach any target),
\item $\displaystyle\sum_{n=0}^{\infty} a_n^2 < \infty$ \quad (the steps are small enough that noise averages out).
\end{enumerate}
\end{definition}

\begin{remark}[Why both conditions are needed]
\label{rem:step-size-intuition}
Condition (i) ensures the iterates can travel an unbounded total distance --- without
it, the algorithm could ``freeze'' before reaching $x^*$ regardless of the initial
condition. Condition (ii) controls the accumulated noise variance: since the noise at
step $n$ is scaled by $a_n$, the total noise variance is proportional to
$\sum a_n^2 \Var(\text{noise})$, which must be finite for convergence. The canonical
choice is $a_n = c / n$ for some $c > 0$ (since $\sum 1/n = \infty$ while
$\sum 1/n^2 < \infty$). More generally, $a_n = c / n^\alpha$ works for any
$\alpha \in (1/2, 1]$.
\end{remark}

\subsection{The iteration}

\begin{definition}[Robbins--Monro iteration]
\label{def:rm-iteration}
Given an initial point $x_0 \in \R^d$, the \emph{Robbins--Monro iteration} is
\begin{equation}
\label{eq:rm}
  x_{n+1} = x_n + a_n H(x_n, Y_n), \quad n \geq 0,
\end{equation}
where $\{a_n\}$ satisfies Definition~\ref{def:step-size} and $Y_0, Y_1, \ldots$ are
the noise samples.
\end{definition}

We decompose the update into its ``signal'' and ``noise'' components. Define the
\emph{martingale difference noise}
\begin{equation}
\label{eq:mds-noise}
  M_{n+1} \coloneqq H(x_n, Y_n) - h(x_n).
\end{equation}
Then the iteration becomes
\begin{equation}
\label{eq:rm-decomposed}
  x_{n+1} = x_n + a_n\bigl[h(x_n) + M_{n+1}\bigr].
\end{equation}
This is the canonical form we shall use throughout: the iterate moves in the direction
$h(x_n)$ (the ``drift'') plus additive noise $M_{n+1}$.

\subsection{Classical convergence theorem}

We state the assumptions precisely and then prove convergence.

\begin{definition}[Standing assumptions for Robbins--Monro]
\label{def:rm-assumptions}
Let $\mathcal{F}_n = \sigma(x_0, Y_0, \ldots, Y_{n-1})$ be the natural filtration.
We assume:
\begin{enumerate}[label=(A\arabic*)]
\item \label{A1} \textbf{Step sizes:} $\{a_n\}$ satisfies Definition~\ref{def:step-size}.
\item \label{A2} \textbf{Martingale difference noise:} $\E[M_{n+1} \mid \mathcal{F}_n] = 0$ a.s.\ for all $n$.
\item \label{A3} \textbf{Bounded noise variance:} There exists $K > 0$ such that
  $\E[\|M_{n+1}\|^2 \mid \mathcal{F}_n] \leq K(1 + \|x_n\|^2)$ a.s.
\item \label{A4} \textbf{Stability (monotonicity):} There exists a unique root
  $x^* \in \R^d$ with $h(x^*) = 0$, and $\langle x - x^*, h(x) \rangle < 0$ for
  all $x \neq x^*$. (That is, $h$ points ``inward'' toward $x^*$.)
\item \label{A5} \textbf{Linear growth:} $\|h(x)\| \leq K(1 + \|x\|)$ for all $x$.
\end{enumerate}
\end{definition}

\begin{theorem}[Robbins--Monro convergence]
\label{thm:rm-convergence}
Under assumptions \ref{A1}--\ref{A5}, the Robbins--Monro iteration~\eqref{eq:rm}
satisfies $x_n \cas x^*$.
\end{theorem}

\begin{proof}
We use the Lyapunov function $V(x) = \|x - x^*\|^2$. Define $e_n = x_n - x^*$.
Then
\begin{align*}
  \|e_{n+1}\|^2
    &= \|e_n + a_n(h(x_n) + M_{n+1})\|^2 \\
    &= \|e_n\|^2 + 2a_n \langle e_n, h(x_n)\rangle
      + 2a_n \langle e_n, M_{n+1}\rangle \\
    &\quad + a_n^2 \|h(x_n) + M_{n+1}\|^2.
\end{align*}
Taking conditional expectation with respect to $\mathcal{F}_n$ and using \ref{A2}:
\begin{align}
  \E[\|e_{n+1}\|^2 \mid \mathcal{F}_n]
    &= \|e_n\|^2 + 2a_n \langle e_n, h(x_n)\rangle
      + a_n^2 \E[\|h(x_n) + M_{n+1}\|^2 \mid \mathcal{F}_n]. \label{eq:rm-lyap}
\end{align}
For the last term, by the inequality $\|a+b\|^2 \leq 2\|a\|^2 + 2\|b\|^2$:
\[
  \E[\|h(x_n) + M_{n+1}\|^2 \mid \mathcal{F}_n]
    \leq 2\|h(x_n)\|^2 + 2\E[\|M_{n+1}\|^2 \mid \mathcal{F}_n].
\]
Using \ref{A3} and \ref{A5}:
\[
  2\|h(x_n)\|^2 + 2\E[\|M_{n+1}\|^2 \mid \mathcal{F}_n]
    \leq 2K^2(1 + \|x_n\|^2) + 2K(1 + \|x_n\|^2)
    \leq C(1 + \|x_n\|^2)
\]
for some constant $C > 0$. Since $\|x_n\|^2 \leq 2\|e_n\|^2 + 2\|x^*\|^2$, we
obtain
\begin{equation}
\label{eq:rm-lyap2}
  \E[\|e_{n+1}\|^2 \mid \mathcal{F}_n]
    \leq \|e_n\|^2 + 2a_n \langle e_n, h(x_n)\rangle + a_n^2 C'(1 + \|e_n\|^2)
\end{equation}
for some $C' > 0$.

By \ref{A4}, $\langle e_n, h(x_n) \rangle < 0$ when $x_n \neq x^*$. Define
$V_n = \|e_n\|^2$ and
\[
  W_n = \sum_{k=0}^{n-1} (-2a_k) \langle e_k, h(x_k)\rangle, \quad
  U_n = \sum_{k=0}^{n-1} a_k^2 C'(1 + \|e_k\|^2).
\]
Then $V_n + W_n - U_n$ is a (nonnegative when adjusted) supermartingale. More
precisely, writing $\beta_n = -2\langle e_n, h(x_n)\rangle > 0$ (for $x_n \neq x^*$),
we have from~\eqref{eq:rm-lyap2}:
\[
  \E[V_{n+1} \mid \mathcal{F}_n] \leq V_n - a_n \beta_n + a_n^2 C'(1 + V_n + 2\|x^*\|^2).
\]

We apply the Robbins--Siegmund theorem (a standard supermartingale convergence result).

\medskip\noindent
\textbf{Robbins--Siegmund Theorem.} \emph{Let $\{V_n\}$, $\{\alpha_n\}$,
$\{\beta_n\}$, $\{u_n\}$ be nonneg\-a\-tive adapted sequences with
$\E[V_{n+1} \mid \mathcal{F}_n] \leq (1+\alpha_n)V_n - \beta_n + u_n$ and
$\sum \alpha_n < \infty$, $\sum u_n < \infty$ a.s. Then $V_n$ converges a.s.\ and
$\sum \beta_n < \infty$ a.s.}

\medskip
In our case, set $\alpha_n = a_n^2 C'$, $u_n = a_n^2 C'(1 + 2\|x^*\|^2)$, and
$\beta_n = a_n(-2\langle e_n, h(x_n)\rangle)$. Since $\sum a_n^2 < \infty$ by
\ref{A1}, we have $\sum \alpha_n < \infty$ and $\sum u_n < \infty$. By
Robbins--Siegmund, $V_n = \|e_n\|^2$ converges a.s.\ (to some finite limit) and
\begin{equation}
\label{eq:beta-summable}
  \sum_{n=0}^{\infty} a_n \bigl(-2\langle e_n, h(x_n)\rangle\bigr) < \infty
  \quad \text{a.s.}
\end{equation}

Suppose for contradiction that $V_n \to L > 0$ with positive probability on some
event $A$. On $A$, $\|e_n\| \to \sqrt{L} > 0$, so $x_n$ stays bounded away from
$x^*$. By the continuity of $h$ and \ref{A4}, there exists $\delta > 0$ such that
$-\langle e_n, h(x_n)\rangle \geq \delta$ for all large $n$ on $A$. Then
\[
  \sum_{n} a_n (-2\langle e_n, h(x_n)\rangle) \geq 2\delta \sum_n a_n = \infty,
\]
contradicting~\eqref{eq:beta-summable}. Therefore $V_n \to 0$ a.s., i.e.,
$x_n \cas x^*$.
\end{proof}

\begin{keyresult}[Robbins--Monro Convergence]
Under a monotonicity condition on $h$ (it points toward the root $x^*$),
martingale difference noise with bounded conditional variance, and the step-size
conditions $\sum a_n = \infty$, $\sum a_n^2 < \infty$, the iteration
$x_{n+1} = x_n + a_n[h(x_n) + M_{n+1}]$ converges almost surely to $x^*$.
\end{keyresult}

\begin{rlconnection}[SGD as Robbins--Monro]
Stochastic gradient descent (SGD) for minimizing $f(\theta) = \E_Z[\ell(\theta, Z)]$
is the Robbins--Monro iteration with $h(\theta) = -\nabla f(\theta)$ and
$H(\theta, Z) = -\nabla_\theta \ell(\theta, Z)$. The update is
$\theta_{n+1} = \theta_n - a_n \nabla_\theta \ell(\theta_n, Z_n)$. Under
convexity of $f$ (which implies the monotonicity condition \ref{A4} for $-\nabla f$),
Theorem~\ref{thm:rm-convergence} guarantees $\theta_n \cas \theta^*$.
\end{rlconnection}


%% ============================================================
%% SECTION 3: THE ODE METHOD
%% ============================================================
\section{The ODE Method}
\label{sec:ode-method}

The Robbins--Monro analysis requires the strong monotonicity condition \ref{A4},
which fails for many RL algorithms. The \emph{ODE method}, developed by Ljung (1977)
and Kushner--Clark (1978), replaces this with a more natural condition: convergence of
the SA iteration follows from the stability of the associated ODE. This is the
primary tool for analyzing RL convergence.

\subsection{The associated ODE}

\begin{definition}[Associated ODE]
\label{def:associated-ode}
Given the SA iteration~\eqref{eq:rm-decomposed}, the \emph{associated ODE} is
\begin{equation}
\label{eq:ode}
  \dot{x}(t) = h(x(t)), \quad x(0) = x_0.
\end{equation}
\end{definition}

The key insight is that for small step sizes, the SA iteration
$x_{n+1} \approx x_n + a_n h(x_n)$ behaves like an Euler discretization of the
ODE~\eqref{eq:ode} with time step $a_n$, plus noise that averages out.

\subsection{Interpolated trajectory}

To make the connection between the discrete SA iteration and the continuous ODE
precise, we construct an interpolated trajectory.

\begin{definition}[Interpolated trajectory]
\label{def:interpolated-traj}
Define the time instants $t_0 = 0$ and $t_n = \sum_{k=0}^{n-1} a_k$ for $n \geq 1$.
Note that $t_n \to \infty$ by condition (i) of Definition~\ref{def:step-size}. The
\emph{interpolated trajectory} $\bar{x} : [0, \infty) \to \R^d$ is the piecewise
constant function
\[
  \bar{x}(t) = x_n \quad \text{for } t \in [t_n, t_{n+1}).
\]
Equivalently, one may use the piecewise linear interpolation; both lead to the same
asymptotic analysis.
\end{definition}

\subsection{Assumptions for the ODE method}

\begin{definition}[ODE method assumptions]
\label{def:ode-assumptions}
We strengthen and modify the assumptions as follows:
\begin{enumerate}[label=(B\arabic*)]
\item \label{B1} \textbf{Step sizes:} As in \ref{A1}.
\item \label{B2} \textbf{Lipschitz drift:} $h : \R^d \to \R^d$ is Lipschitz
  continuous: $\|h(x) - h(y)\| \leq L\|x - y\|$ for all $x, y$ and some $L > 0$.
\item \label{B3} \textbf{Martingale difference noise:}
  $\E[M_{n+1} \mid \mathcal{F}_n] = 0$ a.s.\ and
  $\E[\|M_{n+1}\|^2 \mid \mathcal{F}_n] \leq K(1 + \|x_n\|^2)$ a.s.
\item \label{B4} \textbf{Stability (boundedness):} $\sup_n \|x_n\| < \infty$ a.s.
\item \label{B5} \textbf{Asymptotic stability of the ODE:} The ODE~\eqref{eq:ode}
  has a globally asymptotically stable equilibrium $x^*$ (i.e., $h(x^*) = 0$ and
  $x(t) \to x^*$ as $t \to \infty$ for every initial condition $x(0) \in \R^d$).
\end{enumerate}
\end{definition}

\begin{remark}[Stability assumption \ref{B4}]
\label{rem:stability}
Assumption \ref{B4} --- that the iterates remain bounded --- is the most delicate.
In practice, one either (a) proves it from properties of $h$ (see the Borkar--Meyn
theorem in Section~\ref{subsec:borkar-meyn}), (b) uses projection to a compact set,
or (c) verifies it empirically. For many RL algorithms, boundedness follows from the
structure of the MDP (e.g., bounded rewards and a discount factor $\gamma < 1$).
\end{remark}

\subsection{The noise vanishes asymptotically}

The first step in the ODE method is to show that the accumulated noise becomes
negligible.

\begin{lemma}[Noise term vanishes]
\label{lem:noise-vanishes}
Under \ref{B1}, \ref{B3}, and \ref{B4}, the process
$Z_n = \sum_{k=0}^{n-1} a_k M_{k+1}$ converges a.s. In particular, for any
$T > 0$ and $\varepsilon > 0$,
\[
  \sup_{t_n \leq s \leq t_n + T}
    \Bigl\|\sum_{k=n}^{m(s)-1} a_k M_{k+1}\Bigr\| \to 0 \quad \text{a.s.\ as } n \to \infty,
\]
where $m(s) = \min\{j : t_j > s\}$.
\end{lemma}

\begin{proof}
The sequence $Z_n = \sum_{k=0}^{n-1} a_k M_{k+1}$ is a martingale with respect to
$\mathcal{F}_n$. Its quadratic variation satisfies
\[
  \E[\|Z_n\|^2] = \sum_{k=0}^{n-1} a_k^2 \E[\|M_{k+1}\|^2]
    \leq \sum_{k=0}^{n-1} a_k^2 K(1 + \E[\|x_k\|^2]).
\]
By \ref{B4}, $\|x_k\|$ is bounded a.s., so $\E[\|x_k\|^2] \leq C$ for some $C > 0$.
Thus $\E[\|Z_n\|^2] \leq K(1 + C) \sum_{k=0}^{\infty} a_k^2 < \infty$ by \ref{B1}.
By the martingale convergence theorem, $Z_n$ converges a.s.

For the uniform convergence over windows of length $T$, note that for
$t_n \leq s \leq t_n + T$, the ``tail sum'' $\sum_{k=n}^{m(s)-1} a_k M_{k+1}$
is a portion of the convergent series $\sum a_k M_{k+1}$, so it tends to zero
uniformly over bounded time windows. More precisely, since $Z_n$ converges a.s.,
$\|Z_m - Z_n\| \to 0$ as $n, m \to \infty$, and the number of indices $k$ with
$t_n \leq t_k \leq t_n + T$ increases but $a_k \to 0$, giving the claimed uniform
convergence.
\end{proof}

\subsection{Main convergence theorem}

\begin{theorem}[Convergence via the ODE method]
\label{thm:ode-convergence}
Under assumptions \ref{B1}--\ref{B5}, the SA iteration
$x_{n+1} = x_n + a_n[h(x_n) + M_{n+1}]$ satisfies $x_n \cas x^*$.
\end{theorem}

\begin{proof}
The proof proceeds in three steps.

\medskip
\noindent\textbf{Step 1: Shifted interpolated trajectories.}
For each $n \geq 0$, define the shifted interpolated trajectory
$x^{(n)} : [0, \infty) \to \R^d$ by
\[
  x^{(n)}(0) = x_n, \quad
  x^{(n)}(t) = x_{n+k} \text{ for } t \in [t_{n+k} - t_n, t_{n+k+1} - t_n),
  \quad k \geq 0.
\]
This is the piecewise constant interpolation of the SA iterates starting from time
$t_n$, with the time axis reset to zero.

\medskip
\noindent\textbf{Step 2: Equicontinuity and convergence to ODE solutions.}
By \ref{B4}, the iterates $\{x_n\}$ are bounded a.s. The drift $h$ is Lipschitz
by \ref{B2}, so $\|h(x_n)\| \leq L\|x_n\| + \|h(0)\|$ is bounded. For any
$T > 0$ and $s, t \in [0, T]$ with $s < t$, the trajectory satisfies
\[
  x^{(n)}(t) - x^{(n)}(s) = \sum_{k : t_k \in [s+t_n, t+t_n)} a_k [h(x_k) + M_{k+1}].
\]
The ``signal'' part is bounded by $\|h\|_\infty \cdot (t-s)$ (since
$\sum a_k \approx t - s$ over the corresponding indices). By
Lemma~\ref{lem:noise-vanishes}, the noise part vanishes as $n \to \infty$. Therefore,
the family $\{x^{(n)}\}$ is equicontinuous on $[0, T]$ a.s.\ for large $n$.

By the Arzela--Ascoli theorem, every subsequence of $\{x^{(n)}\}$ has a further
subsequence that converges uniformly on $[0, T]$. Let $x^{(\infty)}$ be any such
limit. We claim $x^{(\infty)}$ satisfies the ODE~\eqref{eq:ode}.

Indeed, the SA iteration gives
\[
  x^{(n)}(t) = x^{(n)}(0) + \int_0^t h(x^{(n)}(s))\,ds + \varepsilon_n(t),
\]
where $\varepsilon_n(t)$ accounts for the discretization error and the noise, both
of which vanish as $n \to \infty$ (the discretization error vanishes because
$a_n \to 0$ by \ref{B1}, and the noise vanishes by Lemma~\ref{lem:noise-vanishes}).
Taking the limit and using the Lipschitz continuity of $h$ to pass through the
integral:
\[
  x^{(\infty)}(t) = x^{(\infty)}(0) + \int_0^t h(x^{(\infty)}(s))\,ds.
\]
Thus $x^{(\infty)}$ is a solution of the ODE~\eqref{eq:ode}.

\medskip
\noindent\textbf{Step 3: Convergence to the equilibrium.}
By \ref{B5}, the equilibrium $x^*$ is globally asymptotically stable for the ODE.
Therefore, for any solution $x(t)$ of the ODE, $x(t) \to x^*$ as $t \to \infty$.

Now we argue that $x_n \to x^*$ a.s. The set of limit points of $\{x_n\}$ is
nonempty (by boundedness \ref{B4}) and compact. From Step~2, any limit point of the
discrete sequence is also an initial condition for an ODE trajectory that must
converge to $x^*$. Moreover, the set of limit points is connected (as the image of
a connected set under a continuous map) and invariant under the ODE flow.

The only compact, connected, ODE-flow-invariant set that is attracted to $x^*$ under
a globally asymptotically stable system is $\{x^*\}$ itself. To see this, suppose
$\bar{x} \neq x^*$ is a limit point. Then there exists a subsequence $x_{n_k} \to
\bar{x}$. The shifted trajectory $x^{(n_k)}$ converges to an ODE solution starting
at $\bar{x}$, which satisfies $x(T) \to x^*$ for large $T$. But then the iterates
after $n_k$ must eventually be close to $x^*$, and by the same argument, they
cannot leave a neighborhood of $x^*$ (by Lyapunov stability of $x^*$). Hence $x^*$
is the unique limit point, and $x_n \cas x^*$.
\end{proof}

\begin{keyresult}[ODE Method -- Main Principle]
A stochastic approximation iteration $x_{n+1} = x_n + a_n[h(x_n) + M_{n+1}]$
with decreasing step sizes converges to the set of asymptotically stable equilibria
of the ODE $\dot{x} = h(x)$, provided the iterates remain bounded and the noise is
a martingale difference sequence.
\end{keyresult}

\subsection{The Borkar--Meyn stability theorem}
\label{subsec:borkar-meyn}

Assumption \ref{B4} (boundedness of iterates) is often the hardest to verify. The
Borkar--Meyn theorem provides sufficient conditions.

\begin{definition}[Borkar--Meyn assumptions]
\label{def:borkar-meyn-assumptions}
In addition to \ref{B1}--\ref{B3}, assume:
\begin{enumerate}[label=(C\arabic*)]
\item \label{C1} $h$ is Lipschitz and the ``scaled'' function
  $h_\infty(x) \coloneqq \lim_{r \to \infty} h(rx)/r$ exists uniformly on compact
  sets.
\item \label{C2} The ODE $\dot{x} = h_\infty(x)$ has the origin as its unique
  globally asymptotically stable equilibrium.
\end{enumerate}
\end{definition}

\begin{theorem}[Borkar--Meyn stability]
\label{thm:borkar-meyn}
Under \ref{B1}--\ref{B3} and \ref{C1}--\ref{C2}, the iterates $\{x_n\}$ are bounded
almost surely: $\sup_n \|x_n\| < \infty$ a.s.
\end{theorem}

\begin{proof}
We give the key steps of the proof.

\medskip
\noindent\textbf{Step 1: Rescaling.}
Suppose, for contradiction, that $\|x_n\| \to \infty$ along a subsequence with
positive probability. Let $r_n = \max(\|x_0\|, \|x_1\|, \ldots, \|x_n\|)$ and
consider the rescaled iterates $\hat{x}_n = x_n / r_n$.

\medskip
\noindent\textbf{Step 2: Rescaled iteration.}
Since $x_{n+1} = x_n + a_n[h(x_n) + M_{n+1}]$, dividing by $r_{n+1}$:
\[
  \hat{x}_{n+1} = \frac{r_n}{r_{n+1}} \hat{x}_n
    + \frac{a_n}{r_{n+1}} [h(r_n \hat{x}_n) + M_{n+1}].
\]
When $r_n \to \infty$, $h(r_n \hat{x}_n) / r_n \to h_\infty(\hat{x}_n)$ by \ref{C1},
and $r_n / r_{n+1} \to 1$ since $a_n \to 0$. The noise term $a_n M_{n+1}/r_{n+1}$
becomes negligible.

\medskip
\noindent\textbf{Step 3: Limiting ODE.}
By the same interpolation argument as in Theorem~\ref{thm:ode-convergence}, the
rescaled trajectory converges (along subsequences) to solutions of the ODE
$\dot{x} = h_\infty(x)$. But $\|\hat{x}_n\| \leq 1$ and $\|\hat{x}_n\| = 1$
infinitely often (when $r_n$ increases). So the limiting trajectory is confined to
the unit ball and visits its boundary.

\medskip
\noindent\textbf{Step 4: Contradiction.}
By \ref{C2}, every solution of $\dot{x} = h_\infty(x)$ converges to the origin.
But the limiting trajectory of $\hat{x}_n$ must repeatedly visit $\|\hat{x}\| = 1$,
contradicting the global asymptotic stability of the origin for the $h_\infty$ ODE.
Therefore $\sup_n \|x_n\| < \infty$ a.s.
\end{proof}

\begin{remark}[Practical significance]
\label{rem:borkar-meyn-practical}
The Borkar--Meyn theorem is crucial because it converts the verification of iterate
boundedness into a condition on the ``behavior of $h$ at infinity.'' For linear SA
where $h(x) = Ax + b$, we have $h_\infty(x) = Ax$, so \ref{C2} reduces to
requiring that $A$ is Hurwitz (all eigenvalues have negative real parts). This is
exactly the condition for TD(0) convergence with linear function approximation.
\end{remark}


%% ============================================================
%% SECTION 4: CONVERGENCE RATE ANALYSIS
%% ============================================================
\section{Convergence Rate Analysis}
\label{sec:convergence-rate}

Having established convergence, we now ask: \emph{how fast} does $x_n \to x^*$?

\subsection{Central limit theorem for SA}

\begin{definition}[Assumptions for the CLT]
\label{def:clt-assumptions}
Assume \ref{B1}--\ref{B5} hold, and additionally:
\begin{enumerate}[label=(D\arabic*)]
\item \label{D1} $h$ is differentiable at $x^*$ with Jacobian
  $A = Dh(x^*) = \nabla h(x^*) \in \R^{d \times d}$.
\item \label{D2} $A$ is Hurwitz: all eigenvalues of $A$ have strictly negative real
  parts. (This is automatic from the asymptotic stability of $x^*$ for the linearized
  ODE.)
\item \label{D3} The step size is $a_n = a / n$ for some constant $a > 0$ such that
  $aA + \frac{1}{2}I$ is Hurwitz (i.e., all eigenvalues of $aA$ have real part less
  than $-1/2$).
\item \label{D4} The conditional covariance converges:
  $\E[M_{n+1} M_{n+1}^T \mid \mathcal{F}_n] \to \Sigma$ in probability as
  $x_n \to x^*$, for some positive definite matrix $\Sigma$.
\end{enumerate}
\end{definition}

\begin{theorem}[CLT for stochastic approximation]
\label{thm:sa-clt}
Under the assumptions of Definition~\ref{def:clt-assumptions}, we have
\[
  \sqrt{n}(x_n - x^*) \cdist \mathcal{N}(0, V),
\]
where $V$ is the unique solution of the \emph{Lyapunov equation}
\begin{equation}
\label{eq:lyapunov}
  (aA + \tfrac{1}{2}I) V + V (aA + \tfrac{1}{2}I)^T + a^2 \Sigma = 0.
\end{equation}
\end{theorem}

\begin{proof}
Define $e_n = x_n - x^*$ and $u_n = \sqrt{n}\, e_n$. Near $x^*$, the drift is
approximately linear:
\[
  h(x_n) = h(x^*) + A e_n + o(\|e_n\|) = A e_n + o(\|e_n\|),
\]
since $h(x^*) = 0$. With $a_n = a/n$, the SA iteration gives
\[
  e_{n+1} = e_n + \frac{a}{n}[A e_n + r_n + M_{n+1}],
\]
where $r_n = h(x_n) - A e_n = o(\|e_n\|)$ is the higher-order remainder.

Substituting $e_n = u_n / \sqrt{n}$:
\[
  \frac{u_{n+1}}{\sqrt{n+1}}
    = \frac{u_n}{\sqrt{n}} + \frac{a}{n}\Bigl[A \frac{u_n}{\sqrt{n}} + r_n + M_{n+1}\Bigr].
\]
Multiplying both sides by $\sqrt{n+1}$ and using $\sqrt{n+1}/\sqrt{n} = 1 + 1/(2n) + O(1/n^2)$:
\begin{align}
  u_{n+1} &= \sqrt{\frac{n+1}{n}} \, u_n + \frac{a\sqrt{n+1}}{n}
    \Bigl[A \frac{u_n}{\sqrt{n}} + r_n + M_{n+1}\Bigr] \notag \\
  &= \Bigl(1 + \frac{1}{2n} + O(1/n^2)\Bigr) u_n
    + \frac{a}{n}\Bigl(1 + O(1/n)\Bigr)[A u_n + \sqrt{n}\, r_n + \sqrt{n}\, M_{n+1}] \notag \\
  &= u_n + \frac{1}{n}\Bigl(aA + \frac{1}{2}I\Bigr) u_n
    + \frac{a}{\sqrt{n}} M_{n+1} + \text{negligible terms}. \label{eq:u-recursion}
\end{align}
The term $\sqrt{n}\, r_n / n = r_n / \sqrt{n}$ vanishes because $r_n = o(\|e_n\|) = o(\|u_n\|/\sqrt{n})$ and $u_n = O(1)$ in the limit.

The recursion~\eqref{eq:u-recursion} is itself a stochastic approximation with step
size $1/n$, drift $g(u) = (aA + \frac{1}{2}I)u$, and diffusion coefficient
$a/\sqrt{n} \cdot \sqrt{n} = a$ times the noise $M_{n+1}$. By \ref{D3}, the matrix
$aA + \frac{1}{2}I$ is Hurwitz, so the origin is a stable equilibrium of the drift.
The process $\{u_n\}$ is a linear stochastic recursion whose stationary distribution
is Gaussian.

Standard results on such recursions (or a direct characteristic function argument)
yield $u_n \cdist \mathcal{N}(0, V)$ where $V$ satisfies the Lyapunov
equation~\eqref{eq:lyapunov}. Specifically, the covariance $V_n = \E[u_n u_n^T]$
satisfies
\[
  V_{n+1} \approx V_n + \frac{1}{n}\Bigl[(aA + \tfrac{1}{2}I) V_n
    + V_n (aA + \tfrac{1}{2}I)^T + a^2 \Sigma\Bigr],
\]
and the fixed point $V_{n+1} = V_n = V$ gives~\eqref{eq:lyapunov}.
\end{proof}

\begin{remark}[Optimal step size]
\label{rem:optimal-step}
The asymptotic covariance $V$ depends on $a$ through the Lyapunov
equation~\eqref{eq:lyapunov}. When $d = 1$ and $A = -\lambda$ with $\lambda > 0$,
the equation becomes $(-2a\lambda + 1)V = a^2 \sigma^2$, giving
$V = a^2 \sigma^2 / (2a\lambda - 1)$. This is minimized when $a = 1/\lambda$,
yielding $V_{\mathrm{opt}} = \sigma^2 / \lambda^2$. In higher dimensions, the
optimal choice requires $a = -(A^T)^{-1}/2$, which depends on unknown quantities.
\end{remark}

\subsection{Polyak--Ruppert averaging}

The optimal step size depends on the Jacobian $A$ at $x^*$, which is typically
unknown. Polyak (1990) and Ruppert (1988) independently proposed a simple device
that achieves the optimal rate \emph{without} knowing $A$: average the iterates.

\begin{theorem}[Polyak--Ruppert averaging]
\label{thm:polyak-ruppert}
Let $\{x_n\}$ be the SA iterates with step size $a_n = a/n^\alpha$ for
$\alpha \in (1/2, 1)$, and define the averaged iterate
\[
  \bar{x}_n = \frac{1}{n} \sum_{k=1}^{n} x_k.
\]
Under assumptions \ref{B1}--\ref{B5} and \ref{D1}--\ref{D4} (with the modified step
size), we have
\[
  \sqrt{n}(\bar{x}_n - x^*) \cdist \mathcal{N}(0, A^{-1} \Sigma (A^{-T})).
\]
This is the optimal asymptotic covariance (in the sense of the Cramer--Rao bound for
the root-finding problem).
\end{theorem}

\begin{proof}
Since $x_n \to x^*$ and $h(x^*) = 0$, we have
\[
  x_{n+1} - x_n = a_n [h(x_n) + M_{n+1}] = a_n [A(x_n - x^*) + r_n + M_{n+1}],
\]
where $r_n = o(\|x_n - x^*\|)$. Summing from $k = 1$ to $n$:
\[
  x_{n+1} - x_1 = A \sum_{k=1}^{n} a_k (x_k - x^*)
    + \sum_{k=1}^{n} a_k r_k + \sum_{k=1}^{n} a_k M_{k+1}.
\]
Rearranging:
\[
  A \sum_{k=1}^{n} a_k (x_k - x^*) = (x_{n+1} - x_1) - \sum_{k=1}^{n} a_k r_k
    - \sum_{k=1}^{n} a_k M_{k+1}.
\]

Define $S_n = \sum_{k=1}^n a_k$. Since $a_k = a/k^\alpha$ with $\alpha < 1$,
$S_n \sim \frac{a}{1-\alpha} n^{1-\alpha}$.

Now, the key observation: $\bar{x}_n - x^* = \frac{1}{n}\sum_{k=1}^n (x_k - x^*)$.
We relate this to the weighted sum above. A careful Abel summation argument shows
that
\[
  \sqrt{n}(\bar{x}_n - x^*)
    = -A^{-1} \cdot \frac{1}{\sqrt{n}} \sum_{k=1}^{n} M_{k+1} + o_P(1).
\]
The dominant term is the sample average of the martingale differences, and the
remainder terms ($x_{n+1} - x_1$ is $o(\sqrt{n})$, and $\sum a_k r_k$ is negligible)
vanish after normalization.

By the martingale CLT, $\frac{1}{\sqrt{n}} \sum_{k=1}^n M_{k+1} \cdist \mathcal{N}(0, \Sigma)$. Therefore
\[
  \sqrt{n}(\bar{x}_n - x^*) \cdist \mathcal{N}(0, A^{-1} \Sigma A^{-T}). \qedhere
\]
\end{proof}

\begin{keyresult}[Polyak--Ruppert Averaging]
By averaging the iterates of a slowly decaying SA scheme ($a_n = O(n^{-\alpha})$,
$\alpha \in (1/2,1)$), the averaged iterate $\bar{x}_n$ achieves the optimal
convergence rate $O(1/\sqrt{n})$ with the minimal asymptotic covariance
$A^{-1}\Sigma A^{-T}$, without knowledge of the Jacobian $A = Dh(x^*)$.
\end{keyresult}

\begin{rlconnection}[Iterate Averaging in RL]
Polyak--Ruppert averaging suggests that RL algorithms should average their iterates
rather than use the last iterate. In practice, \emph{exponential moving averages}
(EMA) are used as a practical approximation: the target network in DQN, the
exponentially averaged value function in some actor-critic methods, and the
``Polyak averaging'' of policy parameters in DDPG/TD3 all implement this principle.
The theoretical guarantee is that averaging provably improves the convergence rate.
\end{rlconnection}


%% ============================================================
%% SECTION 5: TWO-TIMESCALE STOCHASTIC APPROXIMATION
%% ============================================================
\section{Two-Timescale Stochastic Approximation}
\label{sec:two-timescale}

\subsection{Motivation}

Many RL algorithms involve two coupled iterative updates running simultaneously but
at different rates. The canonical example is \emph{actor-critic}: the critic updates
the value function (fast timescale) while the actor updates the policy (slow
timescale). The theory of two-timescale SA, developed by Borkar (1997), provides the
convergence framework for such algorithms.

\subsection{Formulation}

\begin{definition}[Two-timescale SA]
\label{def:two-timescale}
Consider the coupled iterations
\begin{align}
  x_{n+1} &= x_n + a_n [f(x_n, y_n) + M_{n+1}^{(1)}], \label{eq:slow} \\
  y_{n+1} &= y_n + b_n [g(x_n, y_n) + M_{n+1}^{(2)}], \label{eq:fast}
\end{align}
where $x_n \in \R^{d_1}$, $y_n \in \R^{d_2}$, and the step-size sequences satisfy:
\begin{enumerate}[label=(\roman*)]
\item $\sum a_n = \sum b_n = \infty$ and $\sum a_n^2 < \infty$, $\sum b_n^2 < \infty$,
\item $a_n / b_n \to 0$ as $n \to \infty$ \quad (the ``timescale separation'').
\end{enumerate}
The typical choice is $a_n = 1/n$ and $b_n = 1/n^\beta$ for $\beta \in (1/2, 1)$,
or $a_n = 1/n$ and $b_n = 1/n^{2/3}$.
\end{definition}

\begin{remark}[Timescale separation]
\label{rem:timescale}
The condition $a_n / b_n \to 0$ means that the $y$-iteration uses \emph{larger} step
sizes than the $x$-iteration, so $y$ moves faster. From the perspective of the
$y$-update, $x$ changes so slowly that it appears quasi-static. From the perspective
of the $x$-update, $y$ equilibrates so quickly that it appears to have already
converged.
\end{remark}

\subsection{Assumptions}

\begin{definition}[Two-timescale assumptions]
\label{def:two-ts-assumptions}
\begin{enumerate}[label=(E\arabic*)]
\item \label{E1} \textbf{Step sizes:} As in Definition~\ref{def:two-timescale}.
\item \label{E2} \textbf{Lipschitz drifts:} $f : \R^{d_1} \times \R^{d_2} \to \R^{d_1}$
  and $g : \R^{d_1} \times \R^{d_2} \to \R^{d_2}$ are globally Lipschitz.
\item \label{E3} \textbf{Martingale difference noise:}
  $\E[M_{n+1}^{(i)} \mid \mathcal{F}_n] = 0$ and
  $\E[\|M_{n+1}^{(i)}\|^2 \mid \mathcal{F}_n] \leq K(1 + \|x_n\|^2 + \|y_n\|^2)$
  for $i = 1, 2$.
\item \label{E4} \textbf{Stability:} $\sup_n (\|x_n\| + \|y_n\|) < \infty$ a.s.
\item \label{E5} \textbf{Fast timescale equilibrium:} For each fixed $x \in \R^{d_1}$,
  the ODE $\dot{y} = g(x, y)$ has a globally asymptotically stable equilibrium
  $\lambda(x)$, and $\lambda : \R^{d_1} \to \R^{d_2}$ is Lipschitz.
\item \label{E6} \textbf{Slow timescale stability:} The ODE
  $\dot{x} = f(x, \lambda(x))$ has a globally asymptotically stable equilibrium $x^*$.
\end{enumerate}
\end{definition}

\subsection{Convergence theorem}

\begin{theorem}[Two-timescale convergence (Borkar)]
\label{thm:two-timescale}
Under assumptions \ref{E1}--\ref{E6}, the two-timescale SA iterates satisfy
\[
  x_n \cas x^*, \quad y_n \cas \lambda(x^*) \coloneqq y^*.
\]
\end{theorem}

\begin{proof}
The proof is structured in two phases, analyzing each timescale separately.

\medskip
\noindent\textbf{Phase 1: Fast timescale analysis (convergence of $y_n$).}

Consider the $y$-iteration~\eqref{eq:fast}. Since $a_n / b_n \to 0$, the
$x$-iterates change negligibly on the timescale of $b_n$. Formally, on the time
axis defined by $\tau_n = \sum_{k=0}^{n-1} b_k$, the $x$-update has effective
step size $a_n$, which is asymptotically negligible compared to $b_n$.

To make this precise, fix any $T > 0$ and consider the window of indices $[n, m_n]$
where $m_n$ is the largest index such that $\sum_{k=n}^{m_n - 1} b_k \leq T$. On
this window:
\[
  \|x_{m_n} - x_n\| \leq \sum_{k=n}^{m_n-1} a_k \|f(x_k, y_k) + M_{k+1}^{(1)}\|
    \leq C \sum_{k=n}^{m_n-1} a_k,
\]
where $C$ is a bound on $\|f\| + \|M^{(1)}\|$ (which exists by \ref{E4} and
\ref{E2}). Now
\[
  \sum_{k=n}^{m_n-1} a_k = \sum_{k=n}^{m_n-1} \frac{a_k}{b_k} \cdot b_k
    \leq \sup_{k \geq n} \frac{a_k}{b_k} \sum_{k=n}^{m_n-1} b_k
    \leq T \cdot \sup_{k \geq n} \frac{a_k}{b_k} \to 0.
\]
So on the fast timescale, $x$ is indeed quasi-static.

Therefore, the $y$-iteration behaves asymptotically like
\[
  y_{n+1} = y_n + b_n [g(\bar{x}, y_n) + M_{n+1}^{(2)}]
\]
for a ``frozen'' value $\bar{x}$ (which equals the current value of the slowly
drifting $x_n$). By the ODE method (Theorem~\ref{thm:ode-convergence}), this
converges to the equilibrium $\lambda(\bar{x})$ of $\dot{y} = g(\bar{x}, y)$.

More precisely, the $y$-iterates asymptotically track the ODE
$\dot{y} = g(x(t), y)$ where $x(t)$ is the slowly varying $x$-component. Since
$x$ changes slowly, the $y$-component ``quasi-statically'' tracks $\lambda(x_n)$:
for any $\varepsilon > 0$,
\begin{equation}
\label{eq:y-tracking}
  \|y_n - \lambda(x_n)\| < \varepsilon \quad \text{for all large } n, \text{ a.s.}
\end{equation}

\medskip
\noindent\textbf{Phase 2: Slow timescale analysis (convergence of $x_n$).}

Now consider the $x$-iteration~\eqref{eq:slow}. By~\eqref{eq:y-tracking},
$y_n \approx \lambda(x_n)$ for large $n$, so
\[
  f(x_n, y_n) = f(x_n, \lambda(x_n)) + [f(x_n, y_n) - f(x_n, \lambda(x_n))].
\]
The second term is bounded by $L_f \|y_n - \lambda(x_n)\| \to 0$ a.s.\ by
the Lipschitz property of $f$.

Therefore, the $x$-iteration asymptotically behaves like
\[
  x_{n+1} = x_n + a_n [f(x_n, \lambda(x_n)) + \tilde{M}_{n+1}],
\]
where $\tilde{M}_{n+1}$ absorbs the original noise $M_{n+1}^{(1)}$ and the
approximation error from replacing $y_n$ by $\lambda(x_n)$. The associated ODE is
\[
  \dot{x} = f(x, \lambda(x)),
\]
which has $x^*$ as a globally asymptotically stable equilibrium by \ref{E6}.
Applying the ODE method again, $x_n \cas x^*$.

Since $\lambda$ is continuous and $x_n \to x^*$, we conclude
$y_n \to \lambda(x^*) = y^*$ a.s.
\end{proof}

\begin{keyresult}[Two-Timescale Convergence Principle]
In a two-timescale SA system where $a_n/b_n \to 0$:
\begin{enumerate}[leftmargin=2em]
\item The fast iterate $y_n$ ``tracks'' the equilibrium map $\lambda(x_n)$ where
  $g(x, \lambda(x)) = 0$.
\item The slow iterate $x_n$ converges as if $y_n$ were always at equilibrium,
  driven by the averaged dynamics $\dot{x} = f(x, \lambda(x))$.
\item The overall system converges to $(x^*, y^*) = (x^*, \lambda(x^*))$.
\end{enumerate}
\end{keyresult}

\begin{rlconnection}[Actor-Critic as Two-Timescale SA]
In actor-critic methods, the critic updates the value function estimate $w_n$ on
the fast timescale and the actor updates the policy parameters $\theta_n$ on the
slow timescale:
\begin{align*}
  w_{n+1} &= w_n + b_n [\delta_n \phi(s_n) - \text{(critic update)}], \\
  \theta_{n+1} &= \theta_n + a_n [\psi_n - \text{(actor update)}],
\end{align*}
with $a_n / b_n \to 0$. Here $(x_n, y_n) = (\theta_n, w_n)$, $f$ is the policy
gradient direction, and $g$ is the TD error for the current policy. The fast
timescale equilibrium $\lambda(\theta)$ is the true value function
$V^{\pi_\theta}$ for policy $\pi_\theta$. The slow timescale ODE is the policy
gradient flow $\dot{\theta} = \nabla J(\theta)$. Theorem~\ref{thm:two-timescale}
then guarantees convergence to a local optimum of $J(\theta)$ (or a global optimum
if $J$ is concave in $\theta$).
\end{rlconnection}

\subsection{Lock-in probability}

When the ODE has multiple locally stable equilibria (as is typical for nonlinear
systems), the SA iteration converges to one of them. Which one depends on the initial
condition and the noise history.

\begin{proposition}[Lock-in to attractors]
\label{prop:lock-in}
Suppose the ODE $\dot{x} = h(x)$ has finitely many equilibria, each of which is
either asymptotically stable or unstable. Under the ODE method assumptions
\ref{B1}--\ref{B4}, the SA iterates converge a.s.\ to one of the stable equilibria.
The probability of converging to an unstable equilibrium is zero.
\end{proposition}

\begin{proof}
By the ODE method analysis (Theorem~\ref{thm:ode-convergence}, Step 2), the
limit set of $\{x_n\}$ is a compact connected invariant set of the ODE flow.

If this limit set contained an unstable equilibrium $\bar{x}$, then since $\bar{x}$
is unstable, there exists an open neighborhood $U$ of $\bar{x}$ and a direction
along which ODE trajectories leave $U$ exponentially fast. However, the noise in
SA prevents the iterates from remaining exactly on the unstable manifold.

Formally, by the ``avoidance of unstable equilibria'' result (Pemantle, 1990;
Brandiere, 1998): under the martingale noise assumption with nondegenerate
conditional covariance (i.e., $\E[M_{n+1}M_{n+1}^T \mid \mathcal{F}_n]$ is
positive definite), the set of initial conditions $x_0$ for which $x_n$ converges
to an unstable equilibrium has Lebesgue measure zero, and $\Prob(x_n \to \bar{x}) = 0$
for any fixed $x_0$.

The key idea: near an unstable equilibrium, the noise keeps ``kicking'' the iterate
off the unstable manifold, and once it leaves a small neighborhood, the ODE dynamics
carry it toward a stable equilibrium.
\end{proof}

\begin{warning}[Multiple Stable Equilibria in RL]
When the mean-field ODE has multiple stable equilibria, the RL algorithm may converge
to a suboptimal one depending on the initial conditions and exploration noise. This is
a fundamental challenge in policy gradient methods, where the objective $J(\theta)$ is
generally non-concave and has multiple local optima. Two-timescale SA (actor-critic)
inherits this issue. Entropy regularization (as in SAC) and exploration bonuses are
practical techniques to mitigate convergence to poor local optima.
\end{warning}


%% ============================================================
%% SECTION 6: LINEAR STOCHASTIC APPROXIMATION
%% ============================================================
\section{Linear Stochastic Approximation}
\label{sec:linear-sa}

The most important special case for RL is \emph{linear} SA, where the drift function
is affine.

\subsection{Setup}

\begin{definition}[Linear SA]
\label{def:linear-sa}
A \emph{linear stochastic approximation} iteration is
\begin{equation}
\label{eq:linear-sa}
  x_{n+1} = x_n + a_n [A x_n + b + M_{n+1}],
\end{equation}
where $A \in \R^{d \times d}$, $b \in \R^d$, and $h(x) = Ax + b$.
\end{definition}

The associated ODE is $\dot{x} = Ax + b$, a linear ODE with equilibrium
$x^* = -A^{-1}b$ (when $A$ is invertible).

\subsection{Convergence conditions}

\begin{theorem}[Convergence of linear SA]
\label{thm:linear-sa}
Suppose $A$ is Hurwitz (all eigenvalues have strictly negative real parts), the step
sizes satisfy \ref{A1}, the noise satisfies \ref{B3}, and the iterates are bounded
a.s. Then $x_n \cas x^* = -A^{-1}b$.
\end{theorem}

\begin{proof}
The ODE $\dot{x} = Ax + b$ has $x^* = -A^{-1}b$ as its unique equilibrium. Since
$A$ is Hurwitz, the function $V(x) = (x - x^*)^T P (x - x^*)$ is a Lyapunov function
for the ODE, where $P$ is the unique positive definite solution of the Lyapunov matrix
equation $A^T P + P A = -I$. (Such $P$ exists precisely because $A$ is Hurwitz.)

Indeed,
\[
  \dot{V} = 2(x - x^*)^T P (Ax + b)
    = 2(x - x^*)^T P A (x - x^*)
    = (x - x^*)^T (A^T P + PA)(x - x^*) = -\|x - x^*\|^2,
\]
so $V$ is strictly decreasing along nonconstant trajectories. This confirms that
$x^*$ is globally asymptotically stable.

The result then follows directly from the ODE method
(Theorem~\ref{thm:ode-convergence}).
\end{proof}

\begin{remark}[Hurwitz condition for linear SA]
\label{rem:hurwitz-intuition}
The Hurwitz condition on $A$ means that the ``drift'' $h(x) = Ax + b$ always pushes
the iterate toward the equilibrium $x^* = -A^{-1}b$. If any eigenvalue of $A$ had
a positive real part, the ODE would be unstable in the corresponding direction, and
the SA iteration would diverge. For the borderline case (eigenvalues on the imaginary
axis), convergence depends on higher-order terms and the noise structure.
\end{remark}

\subsection{Stability via Borkar--Meyn for linear SA}

\begin{corollary}[Automatic stability for linear SA]
\label{cor:linear-stability}
For linear SA~\eqref{eq:linear-sa} with $A$ Hurwitz, the boundedness assumption
$\sup_n \|x_n\| < \infty$ a.s.\ holds automatically (without needing to assume it).
\end{corollary}

\begin{proof}
Apply the Borkar--Meyn theorem (Theorem~\ref{thm:borkar-meyn}). For linear $h(x) = Ax + b$,
the ``function at infinity'' is $h_\infty(x) = \lim_{r \to \infty} (A(rx) + b)/r = Ax$.
Since $A$ is Hurwitz, the ODE $\dot{x} = Ax$ has the origin as its globally
asymptotically stable equilibrium. Thus \ref{C2} is satisfied, and
Theorem~\ref{thm:borkar-meyn} gives $\sup_n \|x_n\| < \infty$ a.s.
\end{proof}

\begin{theorem}[Complete convergence for linear SA]
\label{thm:linear-sa-complete}
For the linear SA iteration~\eqref{eq:linear-sa} with $A$ Hurwitz, step sizes
satisfying \ref{A1}, and noise satisfying \ref{B3}, we have $x_n \cas -A^{-1}b$
without any additional stability assumptions.
\end{theorem}

\begin{proof}
Combine Corollary~\ref{cor:linear-stability} (which gives boundedness) with
Theorem~\ref{thm:linear-sa} (which gives convergence given boundedness).
\end{proof}

\begin{rlconnection}[TD(0) with Linear Function Approximation]
In TD(0) with linear function approximation $V(s) = \phi(s)^T w$, the parameter
vector $w$ is updated by
\[
  w_{n+1} = w_n + \alpha_n \bigl[R_{n+1} + \gamma \phi(s_{n+1})^T w_n
    - \phi(s_n)^T w_n\bigr] \phi(s_n).
\]
This is linear SA with
\[
  A = \E_\pi[\phi(s)(\gamma \phi(s') - \phi(s))^T], \quad
  b = \E_\pi[R \phi(s)],
\]
where the expectation is over the stationary distribution of the Markov chain
under policy $\pi$. The key question is: \emph{is $A$ Hurwitz?}

Under on-policy sampling (the states are visited according to the stationary
distribution $d^\pi$ of $\pi$), it can be shown that $A$ is negative definite
(hence Hurwitz). Specifically, $A = \Phi^T D(\gamma P^\pi - I)\Phi$ where $D$ is
the diagonal matrix of stationary probabilities, and $\gamma P^\pi - I$ is negative
definite on the range of $\Phi$ (since $\gamma < 1$ and $P^\pi$ has spectral
radius 1).

By Theorem~\ref{thm:linear-sa-complete}, TD(0) with linear function approximation
converges a.s.\ under on-policy sampling. This is the celebrated result of Tsitsiklis
and Van Roy (1997).

\textbf{The Deadly Triad:} Under \emph{off-policy} sampling with function
approximation, $A$ may fail to be Hurwitz, and TD(0) can \emph{diverge}. This is
the theoretical explanation for the deadly triad (function approximation +
bootstrapping + off-policy).
\end{rlconnection}


%% ============================================================
%% SECTION 7: APPLICATIONS TO REINFORCEMENT LEARNING
%% ============================================================
\section{Applications to Reinforcement Learning}
\label{sec:rl-applications}

We now systematically map the general SA theory to specific RL algorithms, verifying
the assumptions of the convergence theorems in each case.

\subsection{TD(0) as stochastic approximation}

\begin{rlconnection}[TD(0) -- Complete SA Analysis]
\textbf{Setting.} Consider a finite MDP with states $\mathcal{S} = \{1, \ldots, S\}$,
policy $\pi$, discount factor $\gamma \in [0, 1)$, transition matrix $P^\pi$, and
reward function $r^\pi(s) = \E[R \mid S = s, \pi]$.

\textbf{The update.} TD(0) updates the value function estimate $V$ by
\[
  V_{n+1}(s_n) = V_n(s_n) + \alpha_n(s_n) \bigl[\underbrace{R_{n+1} + \gamma V_n(s_{n+1}) - V_n(s_n)}_{\delta_n \text{ (TD error)}}\bigr].
\]

\textbf{SA form.} Viewing $V \in \R^S$ as the iterate, define
\[
  h(V) = r^\pi + \gamma P^\pi V - V = (\gamma P^\pi - I)V + r^\pi.
\]
Then $h(V^\pi) = r^\pi + \gamma P^\pi V^\pi - V^\pi = 0$ by the Bellman equation.
The noisy observation at step $n$ (for state $s_n$) is
$H_n = (R_{n+1} + \gamma V_n(s_{n+1}) - V_n(s_n)) e_{s_n}$, where $e_s$ is the
$s$-th standard basis vector.

\textbf{Mean field.} $\E[H_n \mid V_n, s_n] = (r^\pi(s_n) + \gamma (P^\pi V_n)(s_n) - V_n(s_n)) e_{s_n}$, which matches $h(V_n)$ projected onto the $s_n$-coordinate.

\textbf{Associated ODE.} $\dot{V} = (\gamma P^\pi - I)V + r^\pi$. The matrix
$A = \gamma P^\pi - I$ has eigenvalues $\gamma \lambda_i - 1$ where $\lambda_i$ are
eigenvalues of $P^\pi$. Since $|\lambda_i| \leq 1$ (stochastic matrix) and
$\gamma < 1$, all eigenvalues of $A$ satisfy
$\mathrm{Re}(\gamma \lambda_i - 1) \leq \gamma - 1 < 0$. Hence $A$ is Hurwitz.

\textbf{Conclusion.} By Theorem~\ref{thm:linear-sa-complete} (or
Theorem~\ref{thm:ode-convergence}), TD(0) converges a.s.\ to
$V^* = -A^{-1}r^\pi = (I - \gamma P^\pi)^{-1}r^\pi = V^\pi$, the true value
function under $\pi$.

\textbf{Step sizes.} Convergence requires the Robbins--Monro conditions to hold
\emph{per state}: $\sum_n \alpha_n(s) \indicator(s_n = s) = \infty$ and
$\sum_n \alpha_n(s)^2 \indicator(s_n = s) < \infty$ a.s.\ for each $s$.
\end{rlconnection}

\subsection{Q-learning as asynchronous SA}

\begin{rlconnection}[Q-learning -- Asynchronous SA]
\textbf{The update.} Q-learning updates a single state-action entry per step:
\[
  Q_{n+1}(s_n, a_n) = Q_n(s_n, a_n) + \alpha_n(s_n, a_n)
    \bigl[R_{n+1} + \gamma \max_{a'} Q_n(s_{n+1}, a') - Q_n(s_n, a_n)\bigr].
\]

\textbf{SA form.} The mean field $h$ is now \emph{nonlinear} due to the $\max$ operator:
\[
  h(Q)(s,a) = r(s,a) + \gamma \sum_{s'} P(s' \mid s, a) \max_{a'} Q(s', a') - Q(s,a).
\]

\textbf{Associated ODE.} $\dot{Q} = h(Q)$, which is the Bellman optimality operator
$\mathcal{T}^* Q - Q$ in continuous time. Since $\mathcal{T}^*$ is a
$\gamma$-contraction in the sup-norm, the ODE has $Q^*$ (the optimal Q-function) as
its unique globally asymptotically stable equilibrium. To see this, note that the
Lyapunov function $V(Q) = \|Q - Q^*\|_\infty$ satisfies
\[
  \frac{d}{dt}\|Q(t) - Q^*\|_\infty \leq -(1 - \gamma)\|Q(t) - Q^*\|_\infty.
\]

\textbf{Complication: asynchronous updates.} Q-learning updates only one $(s,a)$ pair
per step, making it an \emph{asynchronous} SA scheme. The theory extends via results
of Borkar (1998) and Tsitsiklis (1994): if each $(s,a)$ pair is updated infinitely
often (ensured by an exploring policy), the asynchronous iteration converges to the
same limit as the synchronous ODE.

\textbf{Conclusion.} Under the Robbins--Monro step-size conditions (per state-action
pair) and sufficient exploration, Q-learning converges a.s.\ to $Q^*$.
\end{rlconnection}

\subsection{Policy gradient (REINFORCE) as SA}

\begin{rlconnection}[REINFORCE as SA on the Gradient Ascent ODE]
\textbf{The update.} REINFORCE updates the policy parameters $\theta$ by
\[
  \theta_{n+1} = \theta_n + \alpha_n \widehat{\nabla J(\theta_n)},
\]
where $\widehat{\nabla J(\theta_n)}$ is an unbiased estimate of the policy gradient
$\nabla_\theta J(\theta)$.

\textbf{SA form.} This is Robbins--Monro with $h(\theta) = \nabla J(\theta)$ and
noise $M_{n+1} = \widehat{\nabla J(\theta_n)} - \nabla J(\theta_n)$.

\textbf{Associated ODE.} $\dot{\theta} = \nabla J(\theta)$, which is gradient
ascent on the expected return. Any local maximum $\theta^*$ of $J$ with
$\nabla J(\theta^*) = 0$ and negative definite Hessian is a stable equilibrium.

\textbf{Complication.} The ODE $\dot{\theta} = \nabla J(\theta)$ typically has
multiple stable equilibria (local maxima of $J$). By Proposition~\ref{prop:lock-in},
the SA iterates converge to one of them, but global optimality is not guaranteed.

\textbf{Convergence rate.} With proper step sizes, the CLT
(Theorem~\ref{thm:sa-clt}) applies near each local maximum, giving
$\sqrt{n}(\theta_n - \theta^*) \cdist \mathcal{N}(0, V)$ for an appropriate
covariance $V$. The high variance of REINFORCE estimators means $\Sigma$ (and hence
$V$) is large, motivating the use of baselines.
\end{rlconnection}

\subsection{Actor-critic as two-timescale SA}

\begin{rlconnection}[Actor-Critic -- Two-Timescale Analysis]
\textbf{Setup.} Actor-critic maintains two parameter vectors:
\begin{itemize}[leftmargin=2em]
\item Critic parameters $w \in \R^{d_w}$ (fast timescale, step size $b_n$)
\item Actor parameters $\theta \in \R^{d_\theta}$ (slow timescale, step size $a_n$)
\end{itemize}
with $a_n / b_n \to 0$.

\textbf{Coupled iterations.}
\begin{align*}
  w_{n+1} &= w_n + b_n \bigl[\delta_n \phi(s_n) \bigr], \quad
    \delta_n = R_{n+1} + \gamma \phi(s_{n+1})^T w_n - \phi(s_n)^T w_n, \\
  \theta_{n+1} &= \theta_n + a_n \bigl[\delta_n \nabla_\theta \log \pi_\theta(a_n \mid s_n)\bigr].
\end{align*}

\textbf{Mapping to two-timescale SA.}
\begin{itemize}[leftmargin=2em]
\item $y_n = w_n$, $x_n = \theta_n$.
\item $g(\theta, w) = \E_{\pi_\theta}[\delta^\theta_w \phi(s)]$: the critic drift,
  which is exactly the linear SA drift for TD(0) under policy $\pi_\theta$.
\item $f(\theta, w) = \E_{\pi_\theta}[\delta^\theta_w \nabla_\theta \log \pi_\theta(a \mid s)]$:
  the actor drift, which approximates $\nabla J(\theta)$ when $w$ encodes
  $V^{\pi_\theta}$.
\end{itemize}

\textbf{Verifying assumptions.}
\begin{enumerate}[leftmargin=2em]
\item \ref{E5} (fast equilibrium): For fixed $\theta$, the critic ODE
  $\dot{w} = g(\theta, w)$ is linear SA for TD under $\pi_\theta$. Its equilibrium
  $\lambda(\theta) = w^*(\theta)$ is the TD fixed point (projection of $V^{\pi_\theta}$
  onto the linear space).
\item \ref{E6} (slow stability): With $w = \lambda(\theta)$, the actor ODE becomes
  $\dot{\theta} = \nabla J(\theta)$ (approximately), which ascends the performance
  surface.
\end{enumerate}

\textbf{Conclusion.} By Theorem~\ref{thm:two-timescale}, the critic converges to the
TD fixed point $w^*(\theta^*)$ and the actor converges to a (local) optimum $\theta^*$
of $J(\theta)$.
\end{rlconnection}


%% ============================================================
%% SECTION 8: EXERCISES
%% ============================================================
\section{Exercises}
\label{sec:exercises}

\begin{exercise}[Step-size conditions]
\label{ex:step-sizes}
For each of the following step-size sequences, determine whether both Robbins--Monro
conditions ($\sum a_n = \infty$, $\sum a_n^2 < \infty$) are satisfied:
\begin{enumerate}[label=(\alph*)]
\item $a_n = 1/n$
\item $a_n = 1/n^2$
\item $a_n = 1/\sqrt{n}$
\item $a_n = c/(n + d)$ for constants $c, d > 0$
\item $a_n = 1/(n \log n)$ for $n \geq 2$
\item $a_n = 1/n^{3/4}$
\end{enumerate}
For each case that fails, explain which condition fails and the practical
consequence for an SA algorithm using that step size.
\end{exercise}

\begin{exercise}[Scalar Robbins--Monro]
\label{ex:scalar-rm}
Consider the scalar SA iteration $x_{n+1} = x_n + a_n(c - x_n + \xi_{n+1})$ where
$c \in \R$ and $\{\xi_n\}$ are i.i.d.\ with $\E[\xi_n] = 0$ and
$\E[\xi_n^2] = \sigma^2 < \infty$.
\begin{enumerate}[label=(\alph*)]
\item Identify $h(x)$ and $x^*$.
\item Verify the monotonicity condition \ref{A4}.
\item Show that $A = Dh(x^*) = -1$ (the Jacobian at the equilibrium).
\item With $a_n = a/n$, compute the asymptotic covariance
  $V = \lim_{n \to \infty} n \E[(x_n - x^*)^2]$ using the Lyapunov
  equation~\eqref{eq:lyapunov}.
\item Find the optimal $a$ that minimizes $V$, and compute $V_{\mathrm{opt}}$.
\end{enumerate}
\end{exercise}

\begin{exercise}[ODE for Q-learning]
\label{ex:ode-qlearning}
Consider a two-state, two-action MDP with $\mathcal{S} = \{1, 2\}$,
$\mathcal{A} = \{L, R\}$, $\gamma = 0.9$, deterministic transitions
$P(2 \mid 1, L) = 1$, $P(1 \mid 1, R) = 1$, $P(1 \mid 2, L) = 1$,
$P(2 \mid 2, R) = 1$, and rewards $r(1, L) = 1$, $r(1, R) = 0$,
$r(2, L) = 0$, $r(2, R) = 2$.
\begin{enumerate}[label=(\alph*)]
\item Compute $Q^*(s, a)$ for all $(s, a)$ pairs by solving the Bellman optimality
  equation.
\item Write out the associated ODE $\dot{Q} = h(Q)$ explicitly (as a 4-dimensional
  system).
\item Verify that $Q^*$ is the unique equilibrium.
\item \emph{(Bonus)} Simulate the ODE numerically from several initial conditions
  and verify convergence to $Q^*$.
\end{enumerate}
\end{exercise}

\begin{exercise}[Lyapunov function for contraction]
\label{ex:lyapunov-contraction}
Let $T : \R^d \to \R^d$ be a $\gamma$-contraction in the sup-norm with fixed point
$x^*$: $\|Tx - Tx'\|_\infty \leq \gamma \|x - x'\|_\infty$. Define the ODE
$\dot{x} = Tx - x$.
\begin{enumerate}[label=(\alph*)]
\item Show that $V(x) = \|x - x^*\|_\infty$ is a Lyapunov function satisfying
  $\dot{V} \leq -(1-\gamma) V$.
\item Conclude that $x^*$ is globally asymptotically stable with exponential
  convergence rate $1 - \gamma$.
\item Apply this to the Bellman optimality operator $T = \mathcal{T}^*$ to verify
  the ODE analysis for Q-learning.
\end{enumerate}
\end{exercise}

\begin{exercise}[Two-timescale separation]
\label{ex:two-timescale}
Consider the two-timescale system in $\R^2$:
\begin{align*}
  x_{n+1} &= x_n + a_n(-x_n + y_n + M_{n+1}^{(1)}), \\
  y_{n+1} &= y_n + b_n(-2y_n + x_n + M_{n+1}^{(2)}),
\end{align*}
where $a_n = 1/n$, $b_n = 1/n^{2/3}$, and the noises are i.i.d.\ with zero mean
and unit variance.
\begin{enumerate}[label=(\alph*)]
\item Verify the timescale separation condition $a_n/b_n \to 0$.
\item For the fast timescale: fix $x$ and find the equilibrium $\lambda(x)$
  of $\dot{y} = -2y + x$.
\item For the slow timescale: substitute $y = \lambda(x)$ into the $x$-ODE and
  find the equilibrium $x^*$.
\item Find $(x^*, y^*)$ and verify that the coupled ODE system has this as its
  unique equilibrium.
\item Verify the Hurwitz conditions for both the fast and slow ODEs.
\end{enumerate}
\end{exercise}

\begin{exercise}[Polyak--Ruppert for mean estimation]
\label{ex:polyak-ruppert}
Let $X_1, X_2, \ldots$ be i.i.d.\ with $\E[X_i] = \mu$ and $\Var(X_i) = \sigma^2$.
Consider the SA iteration $\theta_{n+1} = \theta_n + a_n(X_{n+1} - \theta_n)$ with
$a_n = 1/n^{2/3}$.
\begin{enumerate}[label=(\alph*)]
\item Show this is SA with $h(\theta) = \mu - \theta$ and identify $A$ and $\Sigma$.
\item What is the asymptotic distribution of $\sqrt{n}(\theta_n - \mu)$?
  (Use the CLT for SA, Theorem~\ref{thm:sa-clt}, noting the step size is $a/n^\alpha$
  with $\alpha = 2/3$.)
\item Define $\bar{\theta}_n = \frac{1}{n}\sum_{k=1}^n \theta_k$. What is the
  asymptotic distribution of $\sqrt{n}(\bar{\theta}_n - \mu)$?
\item Compare the asymptotic variances from (b) and (c). Which is better?
\item Compare both to the sample mean $\bar{X}_n = \frac{1}{n}\sum_{k=1}^n X_k$.
  What is the asymptotic variance of $\sqrt{n}(\bar{X}_n - \mu)$?
\end{enumerate}
\end{exercise}

\begin{exercise}[Hurwitz verification for TD with linear FA]
\label{ex:hurwitz-td}
Consider a three-state MDP with $\mathcal{S} = \{1, 2, 3\}$, policy $\pi$ inducing
transition matrix
\[
  P^\pi = \begin{pmatrix} 0.2 & 0.5 & 0.3 \\ 0.1 & 0.6 & 0.3 \\ 0.4 & 0.3 & 0.3 \end{pmatrix},
\]
discount $\gamma = 0.95$, and stationary distribution $d^\pi = (0.2, 0.45, 0.35)$.
Let the feature matrix be $\Phi = \begin{pmatrix} 1 & 0 \\ 0 & 1 \\ 1 & 1 \end{pmatrix}$.
\begin{enumerate}[label=(\alph*)]
\item Compute the matrix $A = \Phi^T D(\gamma P^\pi - I)\Phi$ where
  $D = \mathrm{diag}(d^\pi)$.
\item Verify that $A$ is negative definite (hence Hurwitz).
\item Conclude that TD(0) converges with this function approximation scheme.
\end{enumerate}
\end{exercise}

\begin{exercise}[Convergence failure: the deadly triad]
\label{ex:deadly-triad}
Construct a simple two-state MDP with linear function approximation and off-policy
sampling such that the matrix $A$ in the linear SA formulation of TD(0) has an
eigenvalue with positive real part, demonstrating divergence. Specifically:
\begin{enumerate}[label=(\alph*)]
\item Take $\mathcal{S} = \{1, 2\}$, $\gamma = 0.99$, and a behavior policy that
  visits state 1 with probability 0.01 and state 2 with probability 0.99.
\item Let $P^\pi$ be the transition matrix under the target policy $\pi$, with
  $P^\pi(1 \mid 1) = 0.9$, $P^\pi(2 \mid 1) = 0.1$, $P^\pi(1 \mid 2) = 0.1$,
  $P^\pi(2 \mid 2) = 0.9$.
\item Use the feature representation $\phi(1) = 1$, $\phi(2) = 2$ (scalar features).
\item Compute $A = \E_\mu[\phi(s)(\gamma\phi(s') - \phi(s))]$ where $\mu$ is the
  behavior distribution, and show it has a positive eigenvalue.
\item Explain why this SA iteration diverges.
\end{enumerate}
\end{exercise}

\begin{exercise}[SA with Markovian noise]
\label{ex:markov-noise}
The standard SA theory assumes martingale difference noise. In RL, the noise is
typically \emph{Markovian} (since observations come from a Markov chain). Explain how
the ODE method extends to Markovian noise under the following conditions:
\begin{enumerate}[label=(\alph*)]
\item The Markov chain $\{Y_n\}$ is ergodic with stationary distribution $\mu$.
\item $h(x) = \int H(x, y) \, d\mu(y) = \E_\mu[H(x, Y)]$.
\item The ``Poisson equation'' $g(x, y) - Pg(x, y) = H(x, y) - h(x)$ has a
  solution $g$ with suitable regularity.
\end{enumerate}
Sketch how $g$ is used to decompose the noise into a martingale difference term plus
a telescoping (bounded) term, recovering the standard ODE method assumptions.
\end{exercise}

\begin{exercise}[Nonlinear two-timescale: GTD2]
\label{ex:gtd2}
The GTD2 algorithm (Sutton et al., 2009) uses two-timescale SA to perform off-policy
TD learning with linear function approximation. The updates are:
\begin{align*}
  w_{n+1} &= w_n + b_n \bigl[\phi_n(\phi_n - \gamma \phi_n')^T \theta_n - \phi_n \phi_n^T w_n\bigr], \\
  \theta_{n+1} &= \theta_n + a_n \bigl[\delta_n \phi_n - \gamma \phi_n' (\phi_n^T w_n)\bigr],
\end{align*}
where $\phi_n = \phi(s_n)$, $\phi_n' = \phi(s_{n+1})$, $\delta_n = R_{n+1} + \gamma \phi_n'^T \theta_n - \phi_n^T \theta_n$.
\begin{enumerate}[label=(\alph*)]
\item Identify the fast and slow iterates and verify the timescale separation.
\item Compute the fast-timescale equilibrium $\lambda(\theta)$ for the $w$-update.
\item Substitute $\lambda(\theta)$ into the $\theta$-update to obtain the
  slow-timescale ODE.
\item Show that the slow-timescale ODE is gradient descent on the \emph{mean squared
  projected Bellman error} (MSPBE).
\end{enumerate}
\end{exercise}


%% ============================================================
%% SECTION 9: SUMMARY
%% ============================================================
\section{Summary of Key Results}
\label{sec:summary}

\begin{center}
\renewcommand{\arraystretch}{1.4}
\begin{tabular}{@{}p{3.8cm}p{4cm}p{5.5cm}@{}}
\toprule
\textbf{Result} & \textbf{Key Condition} & \textbf{RL Application} \\
\midrule
Robbins--Monro (Thm~\ref{thm:rm-convergence})
  & Monotone $h$, MDS noise, $\sum a_n = \infty$, $\sum a_n^2 < \infty$
  & SGD convergence for convex losses \\
\addlinespace
ODE Method (Thm~\ref{thm:ode-convergence})
  & Bounded iterates, Lipschitz $h$, stable ODE equilibrium
  & Primary tool for all RL convergence proofs \\
\addlinespace
Borkar--Meyn (Thm~\ref{thm:borkar-meyn})
  & $h_\infty$ ODE has stable origin
  & Verifies boundedness for linear SA / TD \\
\addlinespace
SA CLT (Thm~\ref{thm:sa-clt})
  & $a_n = a/n$, Hurwitz Jacobian
  & Rate: $\|x_n - x^*\| = O(1/\sqrt{n})$ \\
\addlinespace
Polyak--Ruppert (Thm~\ref{thm:polyak-ruppert})
  & Average iterates, $a_n = O(n^{-\alpha})$, $\alpha \in (1/2,1)$
  & Optimal rate via iterate averaging (target networks, EMA) \\
\addlinespace
Two-Timescale (Thm~\ref{thm:two-timescale})
  & $a_n/b_n \to 0$, fast/slow ODE stability
  & Actor-critic, GTD, GTD2 convergence \\
\addlinespace
Linear SA (Thm~\ref{thm:linear-sa-complete})
  & $A$ Hurwitz
  & TD(0) with linear FA (on-policy) \\
\addlinespace
Lock-in (Prop~\ref{prop:lock-in})
  & Multiple stable equilibria, nondegenerate noise
  & Policy gradient converges to local optimum, avoids saddle points \\
\bottomrule
\end{tabular}
\end{center}

\bigskip

\begin{center}
\renewcommand{\arraystretch}{1.4}
\begin{tabular}{@{}p{3cm}p{3.5cm}p{3cm}p{3.5cm}@{}}
\toprule
\textbf{RL Algorithm} & \textbf{SA Type} & \textbf{$h(x)$ / Drift} & \textbf{Key SA Theorem} \\
\midrule
TD(0) (tabular) & Async.\ linear SA & $(\gamma P^\pi - I)V + r^\pi$ & ODE Method + Hurwitz \\
\addlinespace
TD(0) (linear FA) & Linear SA & $Aw + b$, $A$ Hurwitz (on-policy) & Thm~\ref{thm:linear-sa-complete} \\
\addlinespace
Q-learning & Async.\ nonlinear SA & $\mathcal{T}^*Q - Q$ ($\gamma$-contraction) & ODE + contraction \\
\addlinespace
REINFORCE & SA (gradient ascent) & $\nabla J(\theta)$ & ODE + lock-in \\
\addlinespace
Actor-Critic & Two-timescale SA & Critic: TD drift; Actor: policy gradient & Thm~\ref{thm:two-timescale} \\
\addlinespace
GTD / GTD2 & Two-timescale SA & Gradient of MSPBE & Thm~\ref{thm:two-timescale} \\
\bottomrule
\end{tabular}
\end{center}

\bigskip
\noindent
The central message of this document is a single principle:

\begin{keyresult}[The SA--ODE Correspondence]
\emph{Every} RL update rule of the form $x_{n+1} = x_n + a_n[H(x_n, Y_n)]$ can be
analyzed by:
\begin{enumerate}[leftmargin=2em]
\item Identifying the mean field $h(x) = \E[H(x, Y)]$.
\item Writing the associated ODE $\dot{x} = h(x)$.
\item Analyzing the stability of the ODE (Lyapunov functions, contraction,
  Hurwitz matrices).
\item Applying the appropriate SA convergence theorem.
\end{enumerate}
If the ODE has a stable equilibrium and the noise/step-size conditions are met, the
RL algorithm converges. If the ODE is unstable (e.g., off-policy TD with FA), the RL
algorithm diverges.
\end{keyresult}

\end{document}
